\documentclass[12pt]{article}
\usepackage[T1]{fontenc}
\usepackage[utf8]{inputenc}
\usepackage{lmodern}
\usepackage{textcomp}
\usepackage{enumitem}
\usepackage{geometry}
\geometry{margin=1in}
\begin{document}
\begin{center}
Answer Key - Astronomy - Graduate Level Assessment 12 \
\small (Answers are ordered exactly as in the question paper.)
\end{center}

\section*{Multiple Choice}
\begin{enumerate}[label=\arabic*.]
\item \textbf{Answer:} D
\\ \textit{Options:} A) Neptune; B) Uranus; C) Saturn; D) all of the above
\\ \textit{Note:} All of the jovian planets-Jupiter, Saturn, Uranus, and Neptune-possess ring systems, although they vary in size, composition, and visibility, making "all of the above" the correct answer.
\item \textbf{Answer:} A
\\ \textit{Options:} A) Hydrogen; B) Helium; C) Carbon; D) Oxygen
\\ \textit{Note:} Hydrogen is the most abundant element in the atmospheres of the Jovian planets-Jupiter, Saturn, Uranus, and Neptune-due to their formation from the primordial solar nebula, which was rich in lighter elements.
\item \textbf{Answer:} D
\\ \textit{Options:} A) They are the shattered remains of collisions between planets.; B) They are chunks of rock or ice that condensed long after the planets and moons had formed.; C) They are chunks of rock or ice that were expelled from planets by volcanoes.; D) They are leftover planetesimals that never accreted into planets.
\\ \textit{Note:} The correct answer is accurate because the Solar Nebular theory posits that asteroids and comets are remnants of the early solar system's formation process, representing planetesimals that did not undergo further aggregation to form planets.
\item \textbf{Answer:} B
\\ \textit{Options:} A) Tunguska Siberia.; B) Chicxulub Crater Yucatan Peninsula in Mexico.; C) Quebec Canada.; D) Meteor Crater in Arizona.
\\ \textit{Note:} The Chicxulub Crater, located on the Yucatan Peninsula in Mexico, is widely accepted as the impact site of a 10-kilometer asteroid that struck Earth approximately 65 million years ago, leading to dramatic environmental changes that contributed to the mass extinction of the dinosaurs.
\item \textbf{Answer:} C
\\ \textit{Options:} A) The nucleus contains most of the atom's mass but almost none of its volume.; B) A neutral atom always has equal numbers of electrons and protons.; C) A neutral atom always has equal numbers of neutrons and protons.; D) The electrons can only orbit at particular energy levels.
\\ \textit{Note:} The statement is incorrect because a neutral atom can have varying numbers of neutrons and protons, as the number of neutrons can differ among isotopes of the same element while maintaining overall electrical neutrality.
\end{enumerate}

\section*{True / False}
\begin{enumerate}[label=\arabic*.]
\setcounter{enumi}{5}
\item \textbf{Answer:} True
\\ \textit{Options:} A) True; B) False
\\ \textit{Note:} The nebular theory suggests that the formation of terrestrial planets occurs in the hotter inner regions of the protoplanetary disk, while jovian planets form in the cooler outer regions, leading to a natural distinction in their numbers and types rather than an equal distribution.
\item \textbf{Answer:} False
\\ \textit{Options:} A) True; B) False
\\ \textit{Note:} The statement is false because Dione is not one of the Galilean moons; it is a moon of Saturn, while the Galilean moons refer specifically to Jupiter's four largest moons: Io, Europa, Ganymede, and Callisto.
\item \textbf{Answer:} True
\\ \textit{Options:} A) True; B) False
\\ \textit{Note:} The statement is true because Mercury's elliptical orbit results in a variation of solar radiation flux, with approximately 9 times the solar energy received at perihelion compared to 4 times at aphelion, reflecting the inverse square law of radiation intensity with respect to distance from the Sun.
\item \textbf{Answer:} True
\\ \textit{Options:} A) True; B) False
\\ \textit{Note:} The statement is true because fresh impact craters are characterized by their ejecta, which are materials expelled during the impact, and raised rims, which form as a result of the explosion and subsequent displacement of the surface material around the crater.
\item \textbf{Answer:} False
\\ \textit{Options:} A) True; B) False
\\ \textit{Note:} The presence of an atmosphere is not the most important factor in determining a planet's volcanic and tectonic history, as internal heat, planetary size, and geological activity also play crucial roles in shaping these processes.
\end{enumerate}

\section*{Long Form Response}
\begin{enumerate}[label=\arabic*.]
\setcounter{enumi}{10}
\item \textbf{Answer:} A parsec is defined as the distance at which one astronomical unit (AU) subtends an angle of one arcsecond from Earth. This unit of measurement is vital in astronomy because it provides a practical means of expressing vast interstellar distances, which can be difficult to comprehend using conventional units like kilometers or miles. The relationship between parsecs and arcseconds arises from the principles of parallax, where the apparent shift of a nearby star against distant background objects, observed from two different positions in Earth's orbit, can be quantified to determine stellar distances. As a result, the parsec serves as a cornerstone in astronomical distance measurement, facilitating the understanding of the scale of the universe and the positioning of celestial bodies relative to Earth.
\\ \textit{Note:} The answer correctly defines a parsec based on its relationship to astronomical units and arcseconds, emphasizing its significance in measuring vast distances in space through the method of parallax.
\item \textbf{Answer:} The nuclei of comets play a significant role in generating meteor showers through a process of gradual disintegration as they approach the Sun. As a comet orbits, solar radiation and gravitational interactions cause the comet's icy nucleus to sublimate, releasing gas and dust particles that form a debris trail along its orbital path. When Earth intersects this trail, it encounters these sand-sized particles, which enter the atmosphere at high velocities, resulting in bright streaks of light known as meteors. This understanding of meteor showers not only highlights the dynamic relationship between celestial bodies but also underscores the significance of cometary activity in shaping our planetary environment. Such insights enhance our comprehension of both the origins of meteor showers and the broader implications for potential impacts on Earth from cometary debris.
\\ \textit{Note:} correct because it accurately describes how the sublimation of a comet's nucleus produces a debris trail that, when intersected by Earth, results in meteor showers due to the high-velocity entry of dust particles into the atmosphere.
\end{enumerate}

\vfill
\noindent\textit{End of Answer Key}
\end{document}